% Source: https://artofproblemsolving.com/community/c3h2024847_formula_for_squares
% Stable identity: aops-topic-2024847-q01
Every such square has a side vector $(p,q)$ with $p>0$ and $q\ge0$: of the four directions in which its sides point, exactly one lies in this quarter-plane. Starting from a suitable vertex $V$, the vertices are
\[
V,\quad V+(p,q),\quad V+(p-q,\,p+q),\quad V+(-q,\,p),
\]
so the square fits exactly in an axis-parallel bounding square of side $k=p+q$, touching each of its sides, with one vertex at distance $q$ along the bottom side from the lower left corner (in the figure, the thick blue side has $(p,q)=(2,1)$).

\begin{asy}
usepackage("amsmath");
size(4cm);
defaultpen(fontsize(10pt));
for(int i=0;i<=3;++i) for(int j=0;j<=3;++j) dot((i,j), gray);
draw((0,0)--(3,0)--(3,3)--(0,3)--cycle, linewidth(0.9));
draw((1,0)--(3,1)--(2,3)--(0,2)--cycle, blue+linewidth(0.9));
draw((2,0)--(3,2)--(1,3)--(0,1)--cycle, red+dashed+linewidth(0.9));
label("$q$", (0.5,0), S); label("$p$", (2,0), S);
draw((1,0)--(3,1), blue+linewidth(1.6));
label("\mbox{}",(min(currentpicture,user=true)+max(currentpicture,user=true))/2,black); // keeps the text after the figure black when LaTeX typesets the labels inline
\end{asy}

Conversely, an axis-parallel $k\times k$ grid square together with an offset $q\in\{0,1,\dots,k-1\}$ determines exactly one inscribed square (offset $0$ is the bounding square itself). The figure shows the three squares with bounding square of side $3$. So each $k\times k$ grid square accounts for exactly $k$ squares, and there are $(n+1-k)^2$ positions for a $k\times k$ grid square. The total is
\[
\sum_{k=1}^{n}k(n+1-k)^2.
\]
Write $m=n+1$. Using $\sum_{k=1}^{m-1}k=\dfrac{(m-1)m}{2}$, $\sum_{k=1}^{m-1} k^2=\dfrac{(m-1)m(2m-1)}{6}$ and $\sum_{k=1}^{m-1} k^3=\dfrac{(m-1)^2m^2}{4}$,
\[
\sum_{k=1}^{m-1}k(m-k)^2=m^2\cdot\dfrac{(m-1)m}{2}-2m\cdot\dfrac{(m-1)m(2m-1)}{6}+\dfrac{(m-1)^2m^2}{4}
=m^2(m-1)\cdot\dfrac{6m-4(2m-1)+3(m-1)}{12}=\dfrac{m^2(m-1)(m+1)}{12}.
\]
So the number of squares is
\[
\boxed{\dfrac{n(n+1)^2(n+2)}{12}}.
\]
For example, $n=1$ gives $1$, $n=2$ gives $6$ (four unit squares, one $2\times2$ square and one tilted square), and $n=3$ gives $20$.
